% 共享排版设置。主文档与补充材料均使用 XeLaTeX。
\usepackage[a4paper,left=23mm,right=23mm,top=22mm,bottom=24mm]{geometry}
\usepackage{amsmath,amssymb,bm}
\usepackage{graphicx}
\usepackage{booktabs,tabularx,array}
\usepackage{enumitem}
\usepackage{xcolor}
\usepackage{caption,subcaption}
\usepackage{listings}
\usepackage{longtable}
\usepackage{makecell}
\usepackage{fancyhdr}
\usepackage[backend=biber,style=numeric,sorting=nyt]{biblatex}
\addbibresource{references.bib}
\usepackage[hidelinks,unicode]{hyperref}

\definecolor{hintgray}{RGB}{115,115,115}
\definecolor{rulegray}{RGB}{210,210,210}
\linespread{1.18}
\setlength{\parindent}{2em}
\setlength{\parskip}{0.15em}
\setlength{\emergencystretch}{2em}
\setlist{nosep,leftmargin=2em,topsep=0.3em}
\renewcommand{\arraystretch}{1.25}
\newcolumntype{Y}{>{\raggedright\arraybackslash}X}
\captionsetup{font=small,labelfont=bf,labelsep=quad}
\ctexset{
  section={format=\large\bfseries,beforeskip=1.2em,afterskip=0.6em},
  subsection={format=\normalsize\bfseries,beforeskip=0.9em,afterskip=0.4em},
  subsubsection={format=\normalsize\bfseries,beforeskip=0.5em,afterskip=0.2em}
}
\pagestyle{fancy}
\fancyhf{}
\fancyhead[L]{\small 安全强化学习方法复现报告}
\fancyhead[R]{\small \DocumentType}
\fancyfoot[C]{\small\thepage}
\setlength{\headheight}{15pt}
\renewcommand{\headrulewidth}{0.3pt}
\lstset{language=Python,basicstyle=\small\ttfamily,
  columns=fullflexible,keepspaces=true,breaklines=true,
  showstringspaces=false,frame=single,rulecolor=\color{rulegray},
  numbers=left,numberstyle=\tiny\color{hintgray},numbersep=8pt}

% 填写后删除相应占位命令。false 仅隐藏提示，不代表内容已经完成。
\newif\ifshowplaceholders
\showplaceholderstrue
\newcommand{\pending}{\ifshowplaceholders\textcolor{hintgray}{待补充}\fi}
\newcommand{\blankcontent}[1][1.2cm]{%
  \par\noindent\ifshowplaceholders{\color{hintgray}\small〔待补充〕}\par\fi
  \vspace{#1}}
\newcommand{\summaryitem}[1]{\par\noindent\textbf{#1}\quad\pending\par\vspace{0.5em}}
% 已填写的概述段落：\summarypara{小标题}{正文}
\newcommand{\summarypara}[2]{\par\noindent\textbf{#1}\quad #2\par\vspace{0.35em}}
% 行内代码/路径：可在 . / _ 处断行
\newcommand{\code}[1]{\nolinkurl{#1}}
\AtBeginDocument{\def\UrlBreaks{\do\.\do\/\do\_\do\-\do\=}}
% 含空格的命令行（\code 会吞掉空格）
\newcommand{\cmd}[1]{\texttt{\detokenize{#1}}}
\newcommand{\reporttitle}[1]{%
  \begin{center}
  {\LARGE\bfseries #1\par}\vspace{0.7em}
  {\normalsize 作者：范金鹏\qquad 日期：2026 年 11 月 8 日\par}
  \end{center}\vspace{0.5em}}

